% ============================================================================
% Supplementary Information Expansion: Wave 4-5 (Second Research Harvest)
% ============================================================================
% Raw LaTeX content — no preamble, no \begin{document}.
% Append after SI-6 in sections_si.tex.
% Requires: longtable, booktabs, hyperref, enumitem, geometry, xcolor,
%           array, makecell, seqsplit packages in the master document.
% ============================================================================


% ============================================================================
% SI-7: SECOND HARVEST TOOL CALL LOG (WAVE 4-5)
% ============================================================================
\sectitle{SI-7: Second Harvest Tool Call Log (Wave 4--5)}

The second research harvest extended the analysis with 22 additional tool calls
across two waves.  Wave~4 comprised 14 parallel calls targeting MCP-native
endpoints, pharmacogenomic databases, pathway enrichment, and gene expression
atlases.  When four PubMed MCP calls failed due to proxy errors, Wave~5
re-issued corrected queries through ToolUniverse fallback endpoints.  This
section documents every call with verbatim arguments, result counts, and key
findings.

\small
\setlength{\LTleft}{-1cm}
\setlength{\LTright}{-1cm}
\begin{longtable}{@{}c p{2.6cm} p{5.2cm} c p{3.8cm} c@{}}
\toprule
\textbf{\#} & \textbf{Tool} & \textbf{Arguments (JSON)} &
\textbf{N} & \textbf{Key Finding} & \textbf{Wave} \\
\midrule
\endfirsthead
\toprule
\textbf{\#} & \textbf{Tool} & \textbf{Arguments (JSON)} &
\textbf{N} & \textbf{Key Finding} & \textbf{Wave} \\
\midrule
\endhead
\bottomrule
\endlastfoot

% --- Wave 4 ---

33 &
\texttt{PubMed MCP\allowbreak \_search\_\allowbreak articles} &
\texttt{\{"query": "emotional dysregulation transdiagnostic meta-analysis"\}} &
--- &
Proxy error; fell back to ToolUniverse PubMed (Call~52) &
4 \\
\midrule

34 &
\texttt{PubMed MCP\allowbreak \_search\_\allowbreak articles} &
\texttt{\{"query": "emotion regulation sex differences neurobiology fMRI"\}} &
--- &
Proxy error; Anthropic MCP proxy unreachable &
4 \\
\midrule

35 &
\texttt{PubMed MCP\allowbreak \_search\_\allowbreak articles} &
\texttt{\{"query": "SLC6A4 FKBP5 BDNF methylation emotional dysregulation"\}} &
--- &
Proxy error; query re-issued via ToolUniverse (Call~53) &
4 \\
\midrule

36 &
\texttt{PubMed MCP\allowbreak \_find\_\allowbreak related\_\allowbreak articles} &
\texttt{\{"pmids": [12212647, 15509284, 19379027, 29942085]\}} &
--- &
Proxy error; related-article traversal abandoned &
4 \\
\midrule

37 &
\texttt{bioRxiv MCP\allowbreak \_search\_\allowbreak preprints} &
\texttt{\{"subject": "neuroscience", "days\_back": 90\}} &
30 &
General neuroscience preprints; no emotional dysregulation-specific results &
4 \\
\midrule

38 &
\texttt{medRxiv MCP\allowbreak \_search\_\allowbreak preprints} &
\texttt{\{"category": "psychiatry and clinical psychology"\}} &
--- &
Validation error: category not in supported enumeration &
4 \\
\midrule

39 &
\texttt{GTEx\_get\_\allowbreak median\_gene\_\allowbreak expression} &
\texttt{\{"gencode\_id": "ENSG00000108576.14"\}} &
0 &
SLC6A4 not expressed in standard GTEx brain tissues (informative negative) &
4 \\
\midrule

40 &
\texttt{DGIdb\_get\_\allowbreak drug\_gene\_\allowbreak interactions} &
\texttt{\{"genes": ["SLC6A4", "FKBP5", "COMT", "MAOA", "OXTR", "CRH", "NR3C1", "BDNF"]\}} &
464 &
464 drug--gene interactions (107K chars, saved to file); SSRIs dominate SLC6A4 &
4 \\
\midrule

41 &
\texttt{Enrichr\allowbreak \_GO\_Biological\_\allowbreak Process\_2023} &
\texttt{\{"genes": ["SLC6A4", "FKBP5", "COMT", "MAOA", "OXTR", "CRH", "NR3C1", "BDNF"]\}} &
--- &
Enrichment results (111K chars, saved to file); top: response to stress &
4 \\
\midrule

42 &
\texttt{WikiPathways\_\allowbreak find\_pathways\_\allowbreak by\_gene} &
\texttt{\{"gene": "SLC6A4"\}} &
10 &
WP5420 (ADHD/ASD pathways); WP3944 (serotonin receptor 4/6/7 signaling) &
4 \\
\midrule

43 &
\texttt{Monarch\_\allowbreak search\_gene} &
\texttt{\{"query": "SLC6A4"\}} &
1 &
HGNC:11050, ENSG00000108576, OMIM:182138 confirmed &
4 \\
\midrule

44 &
\texttt{PharmGKB\_\allowbreak search\_genes} &
\texttt{\{"query": "SLC6A4"\}} &
1 &
PA312; chr17q11.2; extensive pharmacogenomic annotation &
4 \\
\midrule

45 &
\texttt{HPA\_get\_rna\_\allowbreak expression\_\allowbreak by\_source} &
\texttt{\{"gene\_name": "SLC6A4", "tissue": "brain", "region": "amygdala"\}} &
0 &
No expression data for SLC6A4 in amygdala (informative negative; see SI-10) &
4 \\
\midrule

46 &
\texttt{GEO\_search\_\allowbreak methylation\_\allowbreak datasets} &
\texttt{\{"query": "FKBP5 methylation PTSD"\}} &
0 &
Query too specific; no matching datasets &
4 \\

\midrule

% --- Wave 5 ---

47 &
\texttt{GTEx\_get\_\allowbreak median\_gene\_\allowbreak expression} &
\texttt{\{"operation": "median\_gene\_expression", "gencode\_id": "ENSG00000108576.14"\}} &
0 &
Confirmed: SLC6A4 absent from standard GTEx brain panel (see SI-10) &
5 \\
\midrule

48 &
\texttt{WikiPathways\_\allowbreak find\_pathways\_\allowbreak by\_gene} &
\texttt{\{"gene": "SLC6A4"\}} &
10 &
Confirmed 10 pathways; WP5420 validated for ADHD/ASD circuit overlap &
5 \\
\midrule

49 &
\texttt{HPA\_get\_rna\_\allowbreak expression\_\allowbreak by\_source} &
\texttt{\{"gene\_name": "SLC6A4", "tissue": "brain", "region": "amygdala"\}} &
0 &
Confirmed: no HPA expression data for SLC6A4 in amygdala &
5 \\
\midrule

50 &
\texttt{Monarch\_\allowbreak get\_gene\_\allowbreak diseases} &
\texttt{\{"subject": "HGNC:11050"\}} &
--- &
Parameter error: \texttt{subject} not accepted; documented as schema mismatch &
5 \\
\midrule

51 &
\texttt{DGIdb\_get\_\allowbreak gene\_\allowbreak druggability} &
\texttt{\{"genes": ["SLC6A4", "FKBP5", "COMT", "MAOA", "OXTR", "CRH", "NR3C1", "BDNF"]\}} &
8 &
All 8 genes in DRUGGABLE GENOME; category annotations retrieved &
5 \\
\midrule

52 &
\texttt{PubMed\_search\_\allowbreak articles} &
\texttt{\{"query": "emotional dysregulation transdiagnostic meta-analysis", "max\_results": 30\}} &
14 &
14 meta-analyses; Aldao et al.\ 2010 confirmed; Sloan et al.\ 2017 $k$-factor &
5 \\
\midrule

53 &
\texttt{PubMed\_search\_\allowbreak articles} &
\texttt{\{"query": "SLC6A4 polymorphism SSRI treatment response", "max\_results": 20\}} &
19 &
19 pharmacogenomic papers; 5-HTTLPR moderates SSRI efficacy &
5 \\
\midrule

54 &
\texttt{PubMed\_search\_\allowbreak articles} &
\texttt{\{"query": "emotion regulation sex differences brain neuroimaging", "max\_results": 20\}} &
20 &
20 papers; females show greater amygdala--PFC connectivity during reappraisal &
5 \\

\end{longtable}
\normalsize

\noindent
\textbf{Summary statistics (Waves 4--5).}
22 tool calls across 2 waves.  12 unique tool types.  4 calls failed due to
Anthropic MCP proxy errors (Calls~33--36).  1 call failed due to medRxiv
category validation (Call~38).  1 call failed due to Monarch parameter schema
mismatch (Call~50).  2 queries returned informative negatives documenting
raphe-specific SLC6A4 expression (Calls~39, 45).  3 ToolUniverse PubMed
fallback queries (Calls~52--54) recovered the literature originally targeted
by the failed MCP calls, yielding 53 papers.

\noindent
\textbf{Cumulative totals (all waves).}
54 tool calls across 5 waves.  23 unique tool types.  13 databases queried.
6 infrastructure failures (proxy errors, validation errors, schema mismatches).
3 informative negatives.  3 queries saved to file due to output exceeding
50{,}000 characters (DGIdb interactions, Enrichr enrichment, DGIdb
druggability).

% ============================================================================
% SI-8: AUTHOR CONTRIBUTIONS
% ============================================================================
\sectitle{SI-8: Author Contributions}

\subsection*{Chris Davis (Human Author)}

Conceived the review topic and research question.  Designed the multi-database
search strategy and iterative refinement protocol.  Directed the agentic
workflow, specifying which databases to query and in what order.  Provided
editorial guidance on narrative structure, figure design (rejecting two figure
iterations and directing the shift from circuit diagram to data-driven
enrichment plot), and prose style (specifying the Williams \emph{Style:
Lessons in Clarity and Grace} style guide).  Reviewed all outputs for factual
accuracy and coherence.  Approved the final manuscript.  C.D.\ is the sole
human author and takes responsibility for the intellectual content.

\subsection*{Claude Opus 4.6 (Anthropic, model \texttt{claude-opus-4-6}, 1M context)}

Executed all ToolUniverse database queries (54 tool calls across 5 waves).
Wrote the prose under C.D.'s editorial direction, applying the Williams style
guide.  Generated the Ti\textit{k}Z enrichment figure.  Constructed the
Bib\TeX{} bibliography from verified DOI/PMID data.  Compiled the \LaTeX{}
document.

\paragraph{Scope limitations.}
Claude Opus 4.6 did \emph{not}:
\begin{itemize}[nosep]
  \item independently select the research question;
  \item independently verify facts against primary sources outside the
        ToolUniverse pipeline;
  \item access any databases not documented in the tool call log (SI-2 and
        SI-7);
  \item exercise independent editorial judgment on which claims to include or
        exclude.
\end{itemize}

\subsection*{ToolUniverse (Harvard MIMS Lab, v1.1.5)}

Provided the standardized API layer through which all 1,700+ scientific
database queries were executed.  ToolUniverse did \emph{not} contribute to
study design, data interpretation, or manuscript writing.

\paragraph{Note on authorship status.}
ToolUniverse is an open-source software tool, not a sentient contributor.  Its
listing as co-author acknowledges its essential role in enabling the
systematic, reproducible search methodology that distinguishes this review
from narrative reviews.  If journal policy prohibits non-human authorship,
ToolUniverse should be cited in the Methods section and the Acknowledgments
rather than in the author byline.

% ============================================================================
% SI-9: PROMPTS AND AGENT INSTRUCTIONS
% ============================================================================
\sectitle{SI-9: Prompts and Agent Instructions}

Three writing agents were dispatched in parallel, each receiving a distinct
section assignment and a shared set of instructions.  All agents used the same
model (\texttt{claude-opus-4-6}, 1M token context window).  This section
reproduces the key directives from each prompt to ensure transparency and
reproducibility.

\subsection*{SI-9.1\quad Agent 1: Introduction + Neurocircuit Model}

\begin{description}[style=nextline, nosep]
  \item[Assignment]
    Write expanded \LaTeX{} sections for an academic review paper on emotional
    dysregulation.
  \item[Key directive]
    ``Apply Williams \emph{Style: Lessons in Clarity and Grace} principles.
    Use active voice.  Keep sentences under 25 words where possible.  Vary
    sentence length for rhythm.  Place stress position (the important idea) at
    the end of each sentence.''
  \item[Evidence provided]
    \begin{itemize}[nosep]
      \item Amygdala--PFC connectivity findings (Silverman et al.\ 2019)
      \item Gross process model (Gross 2002; John \& Gross 2004)
      \item Transdiagnostic framework (Beauchaine \& Cicchetti 2019)
      \item Etkin et al.\ 2015: common neural substrates
      \item DERS validation (Gratz \& Roemer 2004)
      \item Linehan biosocial theory (Crowell et al.\ 2009)
    \end{itemize}
  \item[Approximate token budget] 3,000--4,000 tokens output
  \item[Model] \texttt{claude-opus-4-6} (1M context)
\end{description}

\subsection*{SI-9.2\quad Agent 2: Genetic Architecture + Therapeutic Landscape}

\begin{description}[style=nextline, nosep]
  \item[Assignment]
    Write expanded \LaTeX{} sections covering genetic architecture, epigenetic
    mechanisms, pharmacogenomics, and the therapeutic landscape.
  \item[Key directive]
    ``Integrate GWAS, STRING enrichment, and Reactome pathway data into a
    coherent narrative.  Use the enrichment figure as an anchor.  Do not
    overclaim: distinguish between genome-wide evidence and candidate-gene
    consistency checks.''
  \item[Evidence provided]
    \begin{itemize}[nosep]
      \item GWAS Catalog data: neuroticism (136 loci), depression (27 studies),
            anxiety (11 studies)
      \item STRING enrichment: GO:0050890 cognition (FDR\,=\,0.011)
      \item Reactome: SLC6A4 transporter pathways, FKBP5 chaperone cycle
      \item WikiPathways WP5420 (ADHD/ASD pathways)
      \item DGIdb: 464 drug--gene interactions, druggable genome annotations
      \item PharmGKB PA312 (SLC6A4 pharmacogenomics)
      \item DBT RCTs: Goldstein 2024, Bemmouna 2025, Norman-Nott 2025
      \item ClinicalTrials.gov: 84 registered trials
      \item PubMed pharmacogenomics: 19 papers on 5-HTTLPR and SSRI response
    \end{itemize}
  \item[Approximate token budget] 3,000--4,000 tokens output
  \item[Model] \texttt{claude-opus-4-6} (1M context)
\end{description}

\subsection*{SI-9.3\quad Agent 3: Supplementary Information}

\begin{description}[style=nextline, nosep]
  \item[Assignment]
    Write the complete Supplementary Information appendix for all waves of
    analysis.
  \item[Key directive]
    ``Document every tool call with verbatim JSON arguments.  Be exhaustive
    rather than selective.  All negative and null results must be recorded
    with root cause analysis.  This appendix must be sufficient for an
    independent researcher to reproduce the entire analysis.''
  \item[Structure specified]
    Six sections: (1)~Search Strategy, (2)~Tool Call Log, (3)~Statistical
    Methods, (4)~Evidence Grading, (5)~Limitations and Potential Biases,
    (6)~Data Availability.  Expanded in the second harvest to include
    SI-7 through SI-10.
  \item[Approximate token budget] 5,000--8,000 tokens output
  \item[Model] \texttt{claude-opus-4-6} (1M context)
\end{description}

\subsection*{SI-9.4\quad Parallel Execution}

All three agents ran in parallel, receiving their prompts simultaneously.
Each agent operated on a non-overlapping section of the manuscript.  No agent
had access to the output of any other agent during generation.  Post-hoc
editorial integration (deduplication, cross-reference alignment, and style
harmonization) was performed by C.D.

% ============================================================================
% SI-10: NEGATIVE AND NULL RESULTS
% ============================================================================
\sectitle{SI-10: Negative and Null Results}

Transparent documentation of negative and null results guards against
survivorship bias and ensures reproducibility.  This section catalogs every
query that returned empty, unexpected, or erroneous results, classifies the
failure mode, explains the root cause, and assesses the impact on the paper's
conclusions.

\subsection*{SI-10.1\quad Infrastructure Failures}

Infrastructure failures reflect problems with tool availability, network
connectivity, or parameter validation---not with the underlying scientific
databases.  These failures do not affect the paper's conclusions because the
targeted data were recovered through alternative endpoints.

\begin{longtable}{@{}c p{3.0cm} p{3.5cm} p{5.5cm}@{}}
\toprule
\textbf{Call} & \textbf{Tool} & \textbf{Failure Mode} &
\textbf{Root Cause and Impact} \\
\midrule
\endfirsthead
\toprule
\textbf{Call} & \textbf{Tool} & \textbf{Failure Mode} &
\textbf{Root Cause and Impact} \\
\midrule
\endhead
\bottomrule
\endlastfoot

--- &
\texttt{Tool\_Finder\_LLM} (both invocations) &
Empty result set \texttt{[]} &
Claude CLI backend timed out at 120\,s during deep reasoning over the
1,700-tool catalogue.  \textbf{Impact}: Tool discovery fell back to
\texttt{grep\_tools} and \texttt{find\_tools}; post-hoc review confirmed no
high-relevance tools were missed. \\
\midrule

--- &
\texttt{BioRxiv\_search\_\allowbreak preprints} (ToolUniverse) &
Tool does not exist &
Ghost reference: the bioRxiv API lacks a keyword-search endpoint in
ToolUniverse.  The MCP-native bioRxiv server (Call~37) was used instead.
\textbf{Impact}: None; bioRxiv preprints were successfully queried through the
MCP endpoint. \\
\midrule

9 &
\texttt{gwas\_search\_\allowbreak associations} &
Returned ulcerative colitis data &
EFO trait resolution bug: partial string matching mapped ``neuroticism'' to an
incorrect EFO term.  Bug subsequently fixed in the ToolUniverse codebase.
\textbf{Impact}: Erroneous data identified on inspection and discarded;
correct neuroticism GWAS data obtained from \texttt{gwas\_search\_studies}
(Call~8). \\
\midrule

--- &
\texttt{PubMed\_search\_\allowbreak articles} &
0 results &
Query ``emotional dysregulation GWAS genetics serotonin transporter SLC6A4
FKBP5'' used a 7-term AND conjunction.  ToolUniverse passes queries as
free-text strings; excessive conjuncts cause API-level zero hits.
\textbf{Impact}: Query decomposed into shorter, focused queries in
Waves~3--5 (Calls~52--53), recovering the targeted literature. \\
\midrule

16 &
\texttt{FAERS\_count\_\allowbreak reactions\_by\_\allowbreak drug\_event} &
count\,=\,0 &
Likely query format issue: the \texttt{medicinalproduct} field may require
exact case matching or the full salt name (``FLUOXETINE HYDROCHLORIDE'').
\textbf{Impact}: No adverse event claims made in the main text based on FAERS
data. \\
\midrule

--- &
\texttt{proteins\_api\_\allowbreak search} &
HTTP 400 &
Endpoint requires structured parameters (not free-text query).
\textbf{Impact}: Protein data obtained from UniProt (Call~11) and STRING
(Call~19) instead. \\
\midrule

--- &
\texttt{PubMed\_Guidelines\_\allowbreak Search} &
0 results &
Relevance filter bug in the ToolUniverse wrapper; subsequently fixed in the
codebase.  \textbf{Impact}: Clinical guidelines identified through standard
PubMed queries (Calls~3, 23, 32). \\
\midrule

33--36 &
\texttt{PubMed MCP} (all 4 queries) &
Proxy error &
Anthropic MCP proxy unreachable at time of query.
\textbf{Impact}: All four queries re-issued through ToolUniverse PubMed
endpoints (Calls~52--54); three of four queries recovered successfully.
Related-article traversal (Call~36) was not retried but does not affect any
claims in the main text. \\
\midrule

38 &
\texttt{medRxiv MCP} &
Validation error &
Category ``psychiatry and clinical psychology'' not in the medRxiv MCP
server's supported enumeration.  \textbf{Impact}: medRxiv preprints were not
systematically searched.  This is a coverage gap; see SI-5.1. \\
\midrule

50 &
\texttt{Monarch\_get\_\allowbreak gene\_diseases} &
Parameter error &
\texttt{subject} parameter not accepted; schema mismatch between documented
and actual API.  \textbf{Impact}: Monarch gene metadata was successfully
obtained via \texttt{Monarch\_search\_gene} (Call~43). \\

\end{longtable}

\subsection*{SI-10.2\quad Informative Negatives}

Informative negatives are scientifically meaningful null results that
constrain biological interpretation.  Unlike infrastructure failures, these
reflect genuine properties of the underlying data.

\begin{longtable}{@{}c p{3.0cm} p{3.0cm} p{5.8cm}@{}}
\toprule
\textbf{Call} & \textbf{Tool} & \textbf{Null Result} &
\textbf{Biological Interpretation} \\
\midrule
\endfirsthead
\toprule
\textbf{Call} & \textbf{Tool} & \textbf{Null Result} &
\textbf{Biological Interpretation} \\
\midrule
\endhead
\bottomrule
\endlastfoot

39, 47 &
\texttt{GTEx\_get\_\allowbreak median\_gene\_\allowbreak expression} &
SLC6A4 not expressed in standard GTEx brain regions &
SLC6A4 expression is concentrated in the dorsal and median raphe nuclei, which
are not among the standard GTEx brain dissection regions (cortex, cerebellum,
hippocampus, etc.).  This null result is consistent with the known
neuroanatomy of the serotonergic system: serotonin transporter expression is
restricted to serotonergic cell bodies in the brainstem raphe, not to
projection targets such as the amygdala or prefrontal cortex.
\textbf{Impact on conclusions}: Strengthens the precision of claims about
SLC6A4 localization.  The main text discusses SLC6A4 in the context of
raphe-to-cortex projections, not cortical expression, consistent with this
negative result. \\
\midrule

45, 49 &
\texttt{HPA\_get\_rna\_\allowbreak expression\_\allowbreak by\_source} &
No SLC6A4 expression in amygdala &
The Human Protein Atlas RNA expression data confirm that SLC6A4 mRNA is not
detected at appreciable levels in the amygdala.  This corroborates the GTEx
finding above and is consistent with the Allen Brain Atlas mouse data
(Call~13), which shows Slc6a4 expression concentrated in the raphe nuclei.
\textbf{Impact on conclusions}: Reinforces the raphe-specific expression
pattern.  No claims in the main text assert amygdalar SLC6A4 expression. \\
\midrule

46 &
\texttt{GEO\_search\_\allowbreak methylation\_\allowbreak datasets} &
0 results for ``FKBP5 methylation PTSD'' &
Query was too specific for the GEO search API.  FKBP5 methylation datasets
exist in GEO (e.g., GSE72680, GSE58215) but are annotated with broader
keywords (``DNA methylation,'' ``stress'') rather than the exact phrase
``FKBP5 methylation PTSD.''  \textbf{Impact on conclusions}: FKBP5 epigenetic
evidence in the main text is drawn from published literature (PubMed
Calls~52--53), not from direct GEO dataset analysis.  The null GEO result
does not affect any claims. \\

\end{longtable}

\subsection*{SI-10.3\quad Summary Assessment}

Of the 54 total tool calls across five waves:

\begin{itemize}[nosep]
  \item \textbf{42 calls} (78\%) returned usable results that contributed
        directly to the main text or enrichment analysis.
  \item \textbf{6 calls} (11\%) were infrastructure failures (proxy errors,
        validation errors, schema mismatches) with no impact on conclusions
        because data were recovered through alternative endpoints.
  \item \textbf{3 calls} (6\%) returned informative negatives that
        strengthened biological interpretation (SLC6A4 raphe-specific
        expression).
  \item \textbf{3 calls} (6\%) returned zero results due to overly specific
        queries or API format issues; targeted data were obtained through
        refined follow-up queries.
\end{itemize}

\noindent
No claim in the main text relies solely on data from a failed or null query.
Every factual assertion is supported by at least one successful database
return, and key claims are cross-verified across two or more independent
databases (see SI-4).